\documentclass[11pt]{article}
\usepackage{amsmath,amssymb,geometry}
\geometry{margin=25mm}
\begin{document}
\section*{Core model: multi-trip VRPTW with an unloading queue}
Source: Section 3.1, PDF pp.3--4, and Section 3.4, PDF p.7,
SPM (44)--(48). This is the full integer trip-set formulation,
over the complete feasible column set.
\paragraph{Physical data.}
$N=\{1,\ldots,n\}$ are collection customers and $0$ is the depot.
Customer $i$ has demand $q_i>0$, service duration $s_i$ and service-start
window $[e_i,l_i]$. Vehicle capacity is $Q$, fleet size is $m$,
and the operating horizon is $[0,T]$. Let $c_{ij}$ be distance and
$t_{ij}$ travel time \emph{including service at the origin}; hence
$t_{ij}=s_i+\text{driving time}_{ij}$ and $s_0=0$.
Depot unloading slots are $u=1,\ldots,P$ of length $L$, starting at
$S_u=(u-1)L$; $S_{P+1}=T$. $C$ is the number of simultaneous unloading
positions, available for simultaneous unloading. Unloading $q$ units
takes $\delta q$ time and reserves complete slots.
\paragraph{Complete trip domain.}
$R$ is the set of all feasible \emph{single} trips. A trip $r$ specifies
an ordered sequence $(0,i_1,\ldots,i_k,0)$ of distinct customers,
a departure time $a_r\ge0$, service-start times $b_{r,i_j}$, a depot
return time $d_r$, and consecutive unloading slots $u_r,\ldots,v_r$.
It is in $R$ precisely when
\begin{align*}
&k\ge1,\quad \sum_{j=1}^kq_{i_j}\le Q,\quad
 b_{r,i_1}\ge a_r+t_{0i_1},\\
&b_{r,i_{j+1}}\ge b_{r,i_j}+t_{i_ji_{j+1}}\quad(j<k),\qquad
 e_{i_j}\le b_{r,i_j}\le l_{i_j}\quad(j=1,\ldots,k),\\
&d_r=b_{r,i_k}+t_{i_k0},\qquad d_r\le S_{u_r},\qquad
 1\le u_r\le v_r\le P,\\
&\delta\sum_{j=1}^kq_{i_j}\le S_{v_r+1}-S_{u_r},\qquad
 f_r:=S_{v_r+1}\le T.
\end{align*}
The unloading reservation consists of consecutive full slots; normally
$v_r-u_r+1=\lceil\delta\sum_jq_{i_j}/L\rceil$. A longer reservation
is dominated, and may be retained without improving the minimum-distance
objective. If the last slot is truncated at $T$, the displayed duration
inequality ensures enough unloading time across the occupied slots. The trip occupies a vehicle continuously on $[a_r,f_r)$,
including travel, queue waiting and the full unloading reservation.
All these quantities are attributes of a column. Trip size is governed by the capacity and timing constraints.
\paragraph{Column coefficients.}
Let $\alpha_{ir}=1$ iff $i$ occurs on $r$,
$\pi_{tr}=\mathbf1\{a_r\le t<f_r\}$ and
$\gamma_{ur}=\mathbf1\{u_r\le u\le v_r\}$.
Define
\[
 c_r=c_{0i_1}+\sum_{j=1}^{k-1}c_{i_ji_{j+1}}+c_{i_k0}.
\]
Thus cost is travel distance. Boundary times use
half-open occupancy so a vehicle can start a new trip exactly when its
previous unloading reservation ends.
\paragraph{One complete integer model.}
\begin{align}
\min_{\theta}\quad &\sum_{r\in R}c_r\theta_r,\tag{44}\\
\text{s.t.}\quad &\sum_{r\in R}\alpha_{ir}\theta_r=1 &&i\in N,\tag{45}\\
&\sum_{r\in R}\pi_{tr}\theta_r\le m &&t\in[0,T],\tag{46}\\
&\sum_{r\in R}\gamma_{ur}\theta_r\le C &&u=1,\ldots,P,\tag{47}\\
&\theta_r\in\{0,1\} &&r\in R.\tag{48}
\end{align}
Selected trip intervals can be assigned to the $m$ identical vehicles
without overlap whenever (46) holds. For a finite column representation,
(46) need only be checked on all event intervals induced by its endpoints.
AF1, AF2, the continuous master and pricing subproblem are excluded.
\end{document}
